\begin{thebibliography}{99}
\bibitem{Apery} R. Ap\'ery,
\emph{Irrationalit\'e de $\zeta(2)$ et $\zeta(3)$},
Ast\'erisque \textbf{61} (1979), 11--13.
\url{https://www.numdam.org/item/AST_1979__61__11_0/}.

\bibitem{BR} K. Ball and T. Rivoal,
\emph{Irrationalit\'e d'une infinit\'e de valeurs de la fonction z\^eta aux entiers impairs},
Invent. Math. \textbf{146} (2001), no.~1, 193--207.
\url{https://doi.org/10.1007/s002220100168}.

\bibitem{Beukers} F. Beukers,
\emph{A note on the irrationality of $\zeta(2)$ and $\zeta(3)$},
Bull. London Math. Soc. \textbf{11} (1979), no.~3, 268--272.
\url{https://doi.org/10.1112/blms/11.3.268}.

\bibitem{BS} M. C. Brambilla and G. Staglian\`o,
\emph{Algebraic boundaries among typical ranks for real binary forms of arbitrary degree},
Found. Comput. Math. \textbf{21} (2021), 1003--1022.
\url{https://doi.org/10.1007/s10208-020-09474-9}.

\bibitem{Lean4} L. de Moura and S. Ullrich,
\emph{The Lean 4 theorem prover and programming language},
in Automated Deduction---CADE 28, Lecture Notes in Computer Science
\textbf{12699}, Springer, 2021, 625--635.
\url{https://doi.org/10.1007/978-3-030-79876-5_37}.

\bibitem{Fischler04} S. Fischler,
\emph{Irrationalit\'e de valeurs de z\^eta (d'apr\`es Ap\'ery, Rivoal, \ldots)},
S\'eminaire Bourbaki, exp.~910, Ast\'erisque \textbf{294} (2004), 27--62.
\url{https://www.numdam.org/item/SB_2002-2003__45__27_0/}.

\bibitem{FSZ} S. Fischler, J. Sprang and W. Zudilin,
\emph{Many odd zeta values are irrational},
Compos. Math. \textbf{155} (2019), no.~5, 938--952.
\url{https://doi.org/10.1112/S0010437X1900722X}.

\bibitem{Arb} F. Johansson,
\emph{Arb: efficient arbitrary-precision midpoint-radius interval arithmetic},
IEEE Trans. Comput. \textbf{66} (2017), no.~8, 1281--1292.
\url{https://doi.org/10.1109/TC.2017.2690633}.

\bibitem{Lai26} L. Lai,
\emph{A note on the number of irrational odd zeta values II},
J. Th\'eor. Nombres Bordeaux \textbf{38} (2026), no.~2, 469--487.
\url{https://doi.org/10.5802/jtnb.1368}.

\bibitem{LY} L. Lai and P. Yu,
\emph{A note on the number of irrational odd zeta values},
Compos. Math. \textbf{156} (2020), no.~8, 1699--1717.
\url{https://doi.org/10.1112/S0010437X20007307}.

\bibitem{Mathlib} The mathlib Community,
\emph{The Lean mathematical library},
in Proceedings of the 9th ACM SIGPLAN International Conference on Certified
Programs and Proofs (CPP 2020), ACM, 2020, 367--381.
\url{https://doi.org/10.1145/3372885.3373824}.

\bibitem{Nesterenko03} Yu. V. Nesterenko,
\emph{Integral identities and constructions of approximations to zeta-values},
J. Th\'eor. Nombres Bordeaux \textbf{15} (2003), no.~2, 535--550.
\url{https://doi.org/10.5802/jtnb.412}.

\bibitem{DLMF} National Institute of Standards and Technology,
\emph{Digital Library of Mathematical Functions},
\S5.11 (Stirling expansions) and \S25.6 (integer arguments of the zeta function).
\url{https://dlmf.nist.gov/5.11}; \url{https://dlmf.nist.gov/25.6}.
Accessed 15 October 2026.

\bibitem{PNT} PrimeNumberTheoremAnd contributors,
\emph{PrimeNumberTheoremAnd}, Lean source development,
revision \path{650d31264be65f4cd6e70c45d8b25d86d482a761}.
\url{https://github.com/AlexKontorovich/PrimeNumberTheoremAnd}.
Exact revision and dependency audit specified in Appendix~\ref{app:reproduction}.

\bibitem{Rivoal00} T. Rivoal,
\emph{La fonction z\^eta de Riemann prend une infinit\'e de valeurs irrationnelles aux entiers impairs},
C. R. Acad. Sci. Paris S\'er. I Math. \textbf{331} (2000), no.~4, 267--270.
\url{https://doi.org/10.1016/S0764-4442(00)01624-4}.

\bibitem{Rivoal02} T. Rivoal,
\emph{Irrationalit\'e d'au moins un des neuf nombres $\zeta(5),\zeta(7),\ldots,\zeta(21)$},
Acta Arith. \textbf{103} (2002), no.~2, 157--167.
\url{https://doi.org/10.4064/aa103-2-5}.

\bibitem{Viola03} C. Viola,
\emph{Birational transformations and values of the Riemann zeta-function},
J. Th\'eor. Nombres Bordeaux \textbf{15} (2003), no.~2, 561--592.
\url{https://doi.org/10.5802/jtnb.414}.

\bibitem{Zudilin01} W. V. Zudilin,
\emph{One of the numbers $\zeta(5),\zeta(7),\zeta(9),\zeta(11)$ is irrational},
Russian Math. Surveys \textbf{56} (2001), no.~4, 774--776.
\url{https://doi.org/10.1070/RM2001v056n04ABEH000427}.

\bibitem{Zudilin04} W. Zudilin,
\emph{Arithmetic of linear forms involving odd zeta values},
J. Th\'eor. Nombres Bordeaux \textbf{16} (2004), no.~1, 251--291.
\url{https://doi.org/10.5802/jtnb.447}.

\bibitem{Zudilin18} W. Zudilin,
\emph{One of the odd zeta values from $\zeta(5)$ to $\zeta(25)$ is irrational.
By elementary means}, SIGMA \textbf{14} (2018), 028, 8~pp.
\url{https://doi.org/10.3842/SIGMA.2018.028}.

\bibitem{Project} \emph{Four-value zeta proof: final contradiction and integration},
Round~29 research archive and final audit, 15 October 2026.
Companion files \path{Four_Value_Lean_Round29.zip} and
\path{Four_Value_Final_Report_2026-10-15.txt}.
Unpublished computational record; contains the formal source, finite data,
verifiers, logs and SHA-256 manifest.
\end{thebibliography}
